\chapter*{The five coupled constructions}
\addcontentsline{toc}{chapter}{The five coupled constructions}

Ordered composition and holomorphic collision ask different questions of a
field algebra. Composition follows the orientation of a transverse interval.
Collision on a smooth curve involves singular coefficients along diagonals
and their normal derivatives. Reconstruction must relate these operations
while preserving both the empty augmentation and the nonempty local vacuum.

The cut product of Chapter~\ref{chap:collision-words} gives a concrete model.
It splits an ordered word into consecutive blocks where the blocks are
disjoint. Positive word length bounds the number of nonempty blocks, so
normalized bar and cobar are finite in each arity. Their inverse equivalence
holds in the enhanced coefficient setting of
Theorem~\ref{thm:words-finite-duality}. The passage from these coefficients to
fields is the comparison problem of Chapter~\ref{cbc:chap:comparisons}.

Five constructions extend this question beyond the word model.
Ordered chiral geometry distinguishes transverse associative composition,
associative composition for chiral convolution, and the chiral Lie bracket.
Each operation has its own bar construction. A mixed distributive law can
relate the two associative directions. Comparing the two orders of their
iterated bars requires that law on the full bisimplicial bar and preservation
of its realizations by the tensor products.

In the finite species model, a mixed coefficient is a direct sum indexed by
pairs of ordered decompositions of a finite set. Projection onto equal
decompositions is an idempotent, the aligned projector. Transporting this
projector to geometric coefficients requires a comparison preserving the
splitting and the transition maps under collision. The finite projection
alone does not construct these data.

Mixed holomorphic--topological Deligne theory begins with the relative
Hochschild complex of a Rees enveloping algebra. Its finite comparison
retains the Rees parameter in the Chevalley differential. Compatible jets
and operadic refinements are additional hypotheses for the complete
open--closed operations.

Igusa--Borcherds theory supplies the root-graded denominator, with its
normalization fixed by the product formula. Calabi--Yau Hall geometry starts
from extension correspondences. Their pullback and pushforward maps define
convolution when the coefficient theory satisfies the required base change.
A quantum double additionally uses a continuous Hopf pairing and the stated
graded completion.

A mixed boundary datum specifies the algebras, their interface, and coherent
actions. Its deformation complex records their simultaneous variation.
A local Batalin--Vilkovisky field theory also requires a local realization
of this datum. An Einstein comparison must identify the physical action
and pairing.

These constructions meet through bar--cobar, open--closed, Hall, current,
trace, and determinant maps. Each comparison must specify its coefficients,
completion, parity, descent, and anomaly hypotheses. The parts that follow
develop the five constructions and the maps between their carriers.
